\section{Introduction}\label{sec:intro}
Atmospheric muons traverse substantial thicknesses of building material. Their attenuation therefore provides information about the integrated material along a trajectory, without placing a radiation source inside the structure. Muon radiography has been applied to large archaeological and geological objects, while detector technology and reconstruction methods have expanded its range of applications \citep{alvarez1970,morishima2017,procureur2018,bonechi2020}. For a building ceiling, the useful signal includes both the broad slab attenuation and the sharper angular shadows of reinforcing beams.

A small tracker samples a limited angular aperture. A single position measures an angular projection rather than a unique three-dimensional density distribution. Moving the same instrument supplies parallax, but two positions still leave a large reconstruction null space. A visually sharp voxel image alone does not establish that its vertical dimensions are measured. An external length or a geometric scale assumption is also necessary: scaling the detector separation and the room together preserves the directions to structural features.

This work examines what can be recovered from two positions of one tracker beneath the cafeteria ceiling at Tel Aviv University. We use a roof exposure of the same detector to define opacity, fit the relative pose by cross-position prediction, and compare a voxel model with height and depth estimators operating on the angular data. A concrete-pinned beam model supplies a separate estimate of vertical beam extent. The on-site beam measurements are excluded from the fitting objective and reported for comparison. The resulting disagreement is part of the result, rather than a criterion for adjusting the analysis.

The contribution is a reproducible comparison of these estimators, their synthetic recovery behavior, and their uncertainty budget. The full-room visualization is treated as a regularized model of the observed opacity, not as a uniquely determined material map.
